\section{Conclusiones}

\begin{enumerate}
	\item El caudal de ventilación requerido para alcanzar 12 renovaciones de aire por hora en el laboratorio de 320 m$^3$ es 64.0 m$^3$/min, equivalente a 3 840 m$^3$/h o 2 260 CFM.
	\item La potencia teórica del ventilador es 0.44 kW para una presión total de 250 Pa y una eficiencia del 60 \%. Se recomienda un motor de 1.0 HP para incluir márgenes de operación y degradación de filtros.
	\item El área de impulsión del ventilador es 0.1333 m$^2$, correspondiente a un diámetro equivalente de 412 mm y un radio de 206.0 mm, con una velocidad de impulsión de 8 m/s.
	\item Las rejillas de exfiltración se dimensionan en tres unidades de 353 mm $\times$ 336 mm cada una, operando a una velocidad facial de 3.0 m/s, lo cual permite mantener la presión positiva sin generar ruido excesivo.
	\item El valor de 12 renovaciones de aire por hora se sustenta en ASHRAE 170, WHO Laboratory Biosafety Manual y NIH DRM, cubriendo el rango superior de laboratorios BSL-2 y el umbral inferior de laboratorios BSL-3.
	\item Los datos de contorno entregados (círculo de inyección de 206.0 mm de radio a 8 m/s y tres rejillas rectangulares de 353 mm $\times$ 336 mm a 3 m/s) son aplicables de manera directa a un modelo CFD para la verificación del campo de presiones y la distribución de velocidades.
\end{enumerate}
